Úvodní poznámka: tato sekce shrnuje praktické nástroje a postupy pro poskytování zpětné vazby na písemné práce v kurzech VUT. Kde to dovolují zdroje, uvedeme konkrétní příklady integrovaných řešení; obecné doporučené postupy jsou označeny jako (general background knowledge).

Krátké rozdělení nástrojů (tierovaný přístup)
\begin{itemize}
  \item Tier 1 — integrované / školní enterprise nástroje: nástroje, které jsou dostupné přímo v rámci školního Office/M365 ekosystému a které lze spravovat centrálně (např. Microsoft 365 Copilot v aplikacích Word, Forms, OneNote apod.). Microsoft 365 Copilot je poskytován jako součást školních Office služeb (ve variantách) a placený doplněk přidává generativní funkce přímo do Word/Excel/PowerPoint/Forms/OneNote atd.; pro školní prostředí existují i webové verze chatbota Copilot Chat \cite{kb_ai_101_tipu_pdf_04e4a115_page_3_1}. 
  \item Tier 2 — specializované asistenční nástroje pro psaní (externí služby): nástroje typu kontrola gramatiky, stylu a čitelnosti (např. Grammarly, Hemingway) a nástroje pro zpětnou vazbu nad strukturou textu. Pozn.: konkrétní popisy funkcí těchto komerčních nástrojů nejsou v dostupných KB výpiscích uvedeny; níže jsou jejich využití a omezení popsány jako (general background knowledge).
  \item Tier 3 — nástroje pro detekci podobnosti a draft‑assistance (originality / průvodní nástroje): nástroje pro kontrolu plagiátorství a práce s návrhy (např. Draft Coach v řešeních typu Turnitin) a platformy, které nabízejí workflow pro revizi draftů. Opět: konkrétní vlastnosti a omezení těchto komerčních produktů jsou zde shrnuty jako (general background knowledge).
\end{itemize}

Doporučené scénáře použití (praktický workflow)
\begin{enumerate}
  \item Nastavte očekávání v sylabu: jasně definujte, které nástroje studenti mohou používat a jak mají deklarovat použití AI (viz vlákno deklarace níže). (general background knowledge)
  \item Požadujte meziverze/drafty: namiřte hodnocení na proces — student odevzdá první draft + krátkou sebereflexi + deklaraci AI; učitel nebo AI‑asistent poskytne strukturovanou formativní zpětnou vazbu; student upraví text a odevzdá finální verzi. Tento postup odpovídá principům milestone‑based workflow používaným v příručce (general background knowledge).
  \item Použijte integrované nástroje tam, kde je potřeba škálovat: pokud máte školní/enterprise Copilot nebo nástroj spravovaný fakultou, využijte jej k hromadnému generování komentářů k formálním aspektům (gramatika, čitelnost, kontrolní otázky) a nechte lektora validovat obsahově klíčové části, zvlášť tam, kde je rozhodující faktická přesnost \cite{kb_ai_101_tipu_pdf_04e4a115_page_3_1}.
  \item Pro matematiku a převod rukopisu: u úloh zahrnujících výpočty využijte nástroje, které dokáží zpracovat postupy (OneNote, math.microsoft.com a funkce „Pokrok v matematice“ byly v praxi popisovány jako užitečné k záznamu postupu a analýze chyb) \cite{kb_ai_101_tipu_pdf_04e4a115_page_8_1}.
  \item Pilotujte a dokumentujte: při prvním cyklu zaznamenejte typické chyby AI‑asistentů a nastavte pravidla retence dat a přístupů (kdo může číst konverzace a drafty). Preferujte školní licence a dohodu o zpracování dat pro nástroje, které budou zpracovávat osobní údaje studentů \cite{kb_ai_101_tipu_pdf_04e4a115_page_3_1}.
\end{enumerate}

Vzorový prompt pro asistenta generujícího formativní zpětnou vazbu
\vspace{0.5em}

\noindent (tento prompt odpovídá doporučenému formátu z knihovny promptů — adaptujte podle rubriky)
\begin{quote}
  Vstupy: (student_submission) (rubric) (declared_AI_use). \\
  Výstup (do ~250 slov): 2 věty chvály, 2 konkrétní návrhy ke zlepšení, 1 návrh na další milník; pokud AI bylo použito, přidej instrukci k ověření specifického tvrzení. Tone: konstruktivní, akční.
\end{quote}
\vspace{0.5em}

Praktický checklist pro učitele před nasazením AI‑podpory zpětné vazby
\begin{itemize}
  \item Je ve sylabu jasně formulováno, jak mají studenti deklarovat použití AI (viz níže uvedená krátká šablona deklarace)? (general background knowledge)
  \item Jsou definovány hranice role AI‑asistenta (např. kontrola jazyka a formy vs. obsahová validace)? (general background knowledge)
  \item Máte nastavené pravidlo sběru postupu pro úlohy, kde je důležitý proces (např. matematické postupy nebo programovací commity)? Pro matematiku zvažte funkce, které umožňují nahrání postupu a jeho analýzu \cite{kb_ai_101_tipu_pdf_04e4a115_page_8_1}.
  \item Je dostupná škálovatelná cesta lidské validace klíčových tvrzení (např. u závěrečných milníků)? (general background knowledge)
  \item Jsou řešeny otázky ochrany osobních údajů a je zajištěna odpovídající licence / DPA pro používané nástroje? Preferujte školní/enterprise varianty tam, kde se zpracovávají studentské artefakty \cite{kb_ai_101_tipu_pdf_04e4a115_page_3_1}.
\end{itemize}

Krátká šablona deklarace použití AI (vložit do instrukcí k odevzdání)
\begin{quote}
  Deklarace použití AI: Uveďte, zda jste použili AI, nástroj a roli (nápověda/editace/generování). Popište ověření správnosti výsledků. Přiložte krátkou sebereflexi (max. 200 slov).
\end{quote}
(vzor převzat z šablon uvedených v příručce; použití doporučeno pro každé odevzdání.)

Poznámky k omezením a rizikům (stručně)
\begin{itemize}
  \item AI nástroje dobře škálují korekturu jazyka a návrhy ke struktuře, avšak obsahovou správnost a hodnocení autorství musí validovat lidský garant. (general background knowledge)
  \item U úloh vyžadujících sledovatelný proces (matematika, programování) preferujte workflow, které sbírá mezi‑verze, snímky postupu, nebo repozitářový záznam (commity). Funkce jako převod rukopisu v OneNote a nástroje pro „pokrok v matematice“ mohou zjednodušit analýzu postupu \cite{kb_ai_101_tipu_pdf_04e4a115_page_8_1}.
  \item Pro demonstrační aktivity a motivaci studentů lze využít výukové světy (např. Minecraft Education) — vhodné zejména pro první seznámení s principy AI v učebním kontextu \cite{kb_ai_101_tipu_pdf_04e4a115_page_15_2}.
\end{itemize}

Závěrem: kombinujte integrované nástroje (kde jsou dostupné a spravované univerzitou) pro škálovatelné aspekty zpětné vazby se strukturovanými lidskými revizemi pro klíčové obsahové části. Pilotní cyklus s jasnými pravidly pro deklaraci AI a sběr postupu je efektivní cestou k udržení akademické integrity a podpory rozvoje studenta. (general background knowledge)